% ============================================================
% kbs_synth.tex — KBS journal: Synthesis / deployment analysis (Section 8).
% Section label: sec:synth. RQ5 (plan): utility per expert-hour as a
% function of coverage; break-even. Requires E1 gain curve + E4 cost ->
% results marked \TODO; the decision-model structure is fixed here.
% Bibliography: kbs_ke_references.bib + paper/references.bib.
% ============================================================

\section{Synthesis: Where Machine Acquisition Pays Off}
\label{sec:synth}

The preceding sections establish three measured functions: acquisition
validity as a function of representation and model
(Section~\ref{sec:valid}), downstream utility as a function of
representation and knowledge-base completeness (Section~\ref{sec:util}),
and acquisition cost in expert time (Section~\ref{sec:cost}). This section
combines them into the deployment analysis that motivates the paper: given
a knowledge base with a known level of subject-knowledge coverage, is
acquiring the missing knowledge by machine worth the expert time required
to verify it, and where is the break-even?

\subsection{The decision model}

We express the decision in the units a knowledge manager controls. Let
$c$ be the current coverage of a knowledge base---the fraction of records
carrying subject knowledge---and let $G(c)$ be the utility gain from
gap-filling at coverage $c$, measured as topical retrieval quality
(Section~\ref{sec:util}: $G(c)$ is decreasing in $c$ and positive
throughout). Let $t_M$ be expert minutes per record to acquire knowledge
manually and $t_V$ the minutes to verify a machine draft
(Section~\ref{sec:cost}), with the verification workflow producing final
knowledge of quality $q_V$ measured against the same gold as manual
acquisition of quality $q_M$ (Section~\ref{sec:valid}). The decision
quantity is \emph{utility per expert-hour}: $G(c) / t_M$ for manual
acquisition against $G(c) / t_V$ for machine acquisition with
verification, and the two are equal at the break-even coverage
$c^{*}$ where the time saving exactly offsets any quality difference.

\subsection{Readings from the measured functions}
\label{sec:synth:readings}

Three structural results follow from the measured functions
\TODO{E1+E4: fill numerical values at journal scale}:

\textbf{1. Machine acquisition pays most where coverage is lowest.}
Because $G(c)$ is monotone decreasing in $c$, the utility per expert-hour
from gap-filling is largest for sparse knowledge bases regardless of the
cost side. This is the operational content of the completeness-dependence
result: the value of the acquisition component is a decreasing function of
how well the base is already covered, so the rational deployment is
backlog and thinly cataloged records, not fully cataloged collections.

\textbf{2. The verification workflow dominates when drafts are near-final.}
If expert verification of a machine draft is faster than manual acquisition
($t_V < t_M$) at equal final quality ($q_V \approx q_M$), the break-even
coverage $c^{*}$ is strictly below 100\%: there is always some coverage
level above which machine-assisted acquisition is the cheaper route to the
same knowledge. The edit-rate measure in Section~\ref{sec:cost} is the
empirical witness for ``drafts are near-final''; where editors change most
drafts substantially, $t_V$ approaches $t_M$ and the advantage narrows
\TODO{E4}.

\textbf{3. Representation targeting concentrates the budget.}
Utility is carried by the thesaurus representation and not by the
hierarchical one (Section~\ref{sec:util}), so $G(c)$ should be computed on
subject-knowledge coverage, and acquisition effort spent on subject
headings---not on classification---when the objective is topical discovery.
Spending acquisition budget on the inert representation buys expert time
with zero utility return.

\subsection{Caveats on the synthesis}

The synthesis inherits the limits of its inputs. The gain curve $G(c)$ is
estimated within the studied collections and should not be assumed to
transfer to collections unlike them (Section~\ref{sec:disc}). The cost
estimates $t_M, t_V$ rest on a small expert panel and are reported
per-cataloger with sensitivity analysis rather than as a single pooled
number \TODO{E4}. And the topical numbers remain upper bounds because the
question distribution is not fully independent of the knowledge being
tested (Section~\ref{sec:disc}). Within these bounds, the synthesis turns
the paper's measurements into an actionable rule: \emph{acquire the missing
subject knowledge by machine, verify it, and spend the saving where
coverage is thinnest}.
